%! AP Statistics — Unit 8, Lesson 8.1 — Homework (Answer Key)
\documentclass[10pt]{article}
\usepackage{apstats-article}
\usepackage{apstats-key}

\begin{document}

\pageheader{Unit 8, Lesson 8.1}{Homework --- Answer Key}
\namedateperiod

\vspace{0.1in}
\textbf{Directions:} For each problem, show your work and explain your reasoning clearly.

\begin{enumerate}

  \item \textbf{Colored Candies.} A bag of 200 colored candies is supposed to contain
        20\% of each of 5 colors. The observed counts are:

        \vspace{0.05in}
        \renewcommand{\arraystretch}{1.3}
        \begin{tabularx}{0.75\linewidth}{Xccccc}
            \textbf{Color}    & Red & Orange & Yellow & Green & Blue \\
            \hline
            \textbf{Observed} & 45  & 38     & 42     & 31    & 44   \\
        \end{tabularx}

        \begin{enumerate}[label=\alph*.]
          \item Find the expected count for each color. Show your calculation.
                \ansline{$0.20 \times 200 = 40$ for each color. Expected: Red 40, Orange 40, Yellow 40, Green 40, Blue 40.}

          \item Which color shows the largest absolute difference
                $|\text{Observed} - \text{Expected}|$ from expected?
                \ansline{Green: $|31 - 40| = 9$. (Red is $|45 - 40| = 5$; Orange $= 2$; Yellow $= 2$; Blue $= 4$.) Green has the largest absolute difference.}

          \item Explain in your own words what question a chi-square goodness-of-fit
                test would answer about this data.
                \ansline{A chi-square GOF test would answer: are the observed differences from 40 (the expected count per color) large enough to conclude that this bag does not follow the claimed 20\%-each distribution, or could these differences plausibly be due to random variation?}
        \end{enumerate}

  \item \textbf{Rolling a Die.} A fair six-sided die is rolled 120 times. You expect
        each face to appear 20 times. The observed counts are:

        \vspace{0.05in}
        \renewcommand{\arraystretch}{1.3}
        \begin{tabularx}{0.75\linewidth}{Xcccccc}
            \textbf{Face}     & 1  & 2  & 3  & 4  & 5  & 6  \\
            \hline
            \textbf{Observed} & 18 & 22 & 19 & 25 & 16 & 20 \\
        \end{tabularx}

        \begin{enumerate}[label=\alph*.]
          \item Compute the difference (Observed $-$ Expected) for each face.
                \ansline{Face 1: $-2$; \; Face 2: $+2$; \; Face 3: $-1$; \; Face 4: $+5$; \; Face 5: $-4$; \; Face 6: $0$.}

          \item Would you be surprised if a fair die produced these results?
                Explain without any formal calculation.
                \ansline{Answers vary. The differences are small relative to an expected count of 20 --- the largest is Face 4 at $+5$. This does not look very surprising; such variation is common with a fair die. A chi-square test would confirm how unusual these gaps are under the assumption of fairness.}
        \end{enumerate}

\end{enumerate}

\begin{reflectionbox}
\textbf{Reflection:}
\ansline{Chi-square tests answer the question: are the differences between what we observed in our data and what we expected (based on a claimed distribution) large enough to conclude that the claimed distribution is wrong, or could those differences be explained by random variation alone?}
\end{reflectionbox}

\end{document}
